\documentclass[11pt]{article}
\usepackage[utf8]{inputenc}
\usepackage{amsmath, amssymb, amsthm, mathtools, bm, geometry, setspace}
\usepackage{booktabs}
\geometry{margin=1in}
\onehalfspacing

\newtheorem{assumption}{Assumption}
\newtheorem{proposition}{Proposition}
\newtheorem{theorem}{Theorem}
\newtheorem{lemma}{Lemma}
\newtheorem{definition}{Definition}
\newtheorem{remark}{Remark}
\newtheorem{corollary}{Corollary}

\title{Measuring Causal Effects under Opaque Targeting:\\
Point and Partial Identification in Overlay Experiments}
\author{Qifan Han}
\date{March 2026}

\begin{document}
\maketitle

\begin{abstract}
Large digital platforms run hundreds of simultaneous A/B tests, deploying proprietary
algorithms to decide which users receive each treatment. From the perspective of a
single advertiser or external researcher, the platform's targeting rule is opaque: the
platform's latent assignment decision is unobserved, and realized treatment receipt
depends on both the experimenter's own randomization and the platform's black-box
logic. Standard experimental designs that randomize eligibility recover only the Average
Treatment Effect on the Treated (ATT) for the platform's latently selected subpopulation,
not the Average Treatment Effect (ATE) that governs population-level investment and
scale-up decisions. This paper develops identification theory for this \emph{overlay
experiment} setting, in which the econometrician randomizes eligibility while the
platform continues to exercise its opaque targeting rule.

We establish two main results. First, when an excluded shifter---a variable that affects
the platform's targeting propensity but not potential outcomes---is available, the
overlay admits a local-IV representation, point-identifying the Marginal Treatment
Effect (MTE) and, under full support, the ATE. Second, when no excluded shifter is
available, we derive sharp bounds on the ATE under a latent-index model with monotone
selection on gains and nonnegative treatment effects: $ATE(x) \in [ITT(x),\, ATT(x)]$.
We introduce the \emph{Selection Ratio} $R(x) = ATT(x)/ATE(x) \geq 1$ as a summary
measure of algorithmic cherry-picking. Our framework is directly applicable to
settings where multiple firms simultaneously overlay experiments on the same platform
and no single firm can observe the others' targeting rules.
\end{abstract}

\section{Introduction}

Large-scale digital platforms---Google, Meta, Amazon, TikTok---run hundreds or thousands
of simultaneous A/B tests every day to iterate their products, algorithms, and
marketing strategies. A single platform may support experiments from dozens of internal
teams and thousands of external advertisers concurrently. As documented by Tang et al.\
(2010) and Kohavi et al.\ (2020), the standard engineering solution to this challenge
is an \emph{orthogonal traffic assignment} system: each experiment is randomized
independently, so that any user may be simultaneously enrolled in multiple experiments
from different teams or advertisers.

This orthogonal design creates a fundamental asymmetry between the platform and any
individual experimenter. The platform sees the joint assignment vector across all
experiments and can---in principle---identify interactions between treatments (Ye et al.\
2025). But an individual advertiser or external researcher sees only their own
randomization. From their perspective, the platform's allocation of their specific
treatment to eligible users is shaped by a proprietary optimization algorithm that
aggregates bids, targeting rules, and signals from all other concurrent campaigns. This
algorithm is \emph{opaque}: the experimenter cannot observe whom the platform would
have chosen to treat absent their eligibility restriction, let alone recover the full
joint assignment distribution.

This paper formalizes the identification problem that arises in this environment.
An experimenter randomizes an eligibility variable $Z_i \in \{0,1\}$---for instance,
a random subset of users is made eligible to receive a sponsored post, a recommendation
override, or a coupon. The platform then decides, within the eligible set, whom to
actually treat, using its own latent targeting rule. Let $D_i \in \{0,1\}$ denote the
platform's latent intent to treat user $i$. Realized treatment receipt is
\[
  W_i = D_i Z_i,
\]
reflecting the fact that treatment occurs only when both the platform intends to treat
($D_i = 1$) and the experimenter has made the user eligible ($Z_i = 1$). When $Z_i = 0$,
the platform's intent is irrelevant---the user is blocked regardless.

This structure, which we call an \emph{overlay experiment}, arises naturally in several
practically important settings:
\begin{itemize}
  \item \textbf{Multi-experiment platforms (within-firm):} A firm's marketing team
    randomly activates a coupon overlay on top of the recommendation team's active
    personalization algorithm. Neither team sees the other's full targeting rule. Each
    team's realized treatment is $W_i = D_i^{\text{other}} \cdot Z_i^{\text{own}}$.
  \item \textbf{Multi-firm advertising:} Multiple advertisers simultaneously bid for the
    same users on a platform. Advertiser A randomizes campaign eligibility, but the
    platform's actual delivery depends on Advertiser A's bids, Advertiser B's
    concurrent campaign, and the platform's auction mechanism---all unobservable to A.
  \item \textbf{Third-party auditing:} An external researcher or regulator cannot force
    a platform to run a clean experiment. They can only randomize incentives or
    eligibility for a sample of users and observe outcomes, while the platform's
    targeting algorithm continues to operate.
\end{itemize}

In all three cases, the experimenter faces a one-sided noncompliance problem: when
$Z_i = 1$, treatment receipt is $W_i = D_i$, which the experimenter can observe; but
the platform's selection rule governing $D_i$ is unknown. As a consequence, as we
formalize in Proposition~\ref{prop:att} below, the passive overlay identifies only the
\emph{average treatment effect on the platform-treated} (ATT)---the causal effect for
the users the platform would have selected. In general, $ATT(x) \neq ATE(x)$, because
the platform endogenously selects users with larger gains from treatment.

A recent and important paper by Ye et al.\ (2025) addresses a related but distinct
problem: when the \emph{platform itself} runs $m$ simultaneous A/B tests with an
\emph{orthogonal, known} assignment mechanism, how can it estimate ATEs for all $2^m$
treatment combinations from only $O(m)$ observed combinations? Their Debiased Deep
Learning (DeDL) framework solves this interpolation problem using DML and structured
neural networks. Crucially, their approach requires that the treatment assignment
mechanism $v(\mathbf{t} \mid x) = \mathbb{P}[\mathbf{T} = \mathbf{t} \mid X = x]$ is
\emph{known}---it is the platform's own randomization rule. Our paper addresses the
complementary case where an external experimenter faces an \emph{unknown} propensity
score $p(x) = \mathbb{P}(D_i = 1 \mid X_i = x)$. The identification strategy must
therefore be entirely different: rather than interpolating across known treatment
combinations, we exploit (i)~an excluded shifter to obtain a local-IV representation,
or (ii)~shape restrictions on the MTE schedule to construct informative bounds.

\textbf{Main contributions.} This paper makes three contributions.

\textbf{1. ATT identification from a passive overlay (Proposition~\ref{prop:att}).}
We show that the ratio of the intent-to-treat effect to the treatment rate,
$ITT(x)/p(x)$, identifies the ATT. This is the \emph{maximal} object identified
from a passive overlay without additional structure, regardless of how the
platform targets users.

\textbf{2. Point identification of the MTE and ATE with an excluded shifter
(Theorem~\ref{thm:pointid}).}
When a variable $S_i$ shifts the platform's propensity score without directly
affecting potential outcomes, the overlay admits a local-IV representation:
$MTE(x, p) = \partial_p \, E[Y_i \mid Z_i = 1, X_i = x, P_i = p]$,
where $P_i = \pi(X_i, S_i)$ is the platform propensity score. Under full support
of $P_i$, the ATE is identified. This section bridges our setting to the
Heckman--Vytlacil (2005) MTE framework.

\textbf{3. Sharp partial identification bounds on the ATE
(Theorem~\ref{thm:sharp}).}
Without an excluded shifter, we derive sharp bounds under a latent-index model
with (i)~monotone selection on gains and (ii)~nonnegative treatment effects:
\[
  ATE(x) \in [ITT(x),\, ATT(x)].
\]
The bounds are sharp: both endpoints are attained by admissible data-generating
processes. We introduce the \emph{Selection Ratio} $R(x) = ATT(x)/ATE(x)$
as a practical summary of the platform's algorithmic cherry-picking efficiency.

\textbf{Paper organization.} Section~\ref{sec:setup} introduces the experimental setup
and notation. Section~\ref{sec:passive} characterizes what a passive overlay identifies.
Section~\ref{sec:pointid} develops point identification with an excluded shifter.
Section~\ref{sec:partial} derives sharp bounds without an excluded shifter and
introduces the Selection Ratio. Section~\ref{sec:lit} reviews related literature.
Section~\ref{sec:proofs} contains all proofs.

\section{Related Literature}
\label{sec:lit}

\textbf{Marginal treatment effects.}
Our point identification result in Section~\ref{sec:pointid} is an application of the
MTE framework developed by Heckman and Vytlacil (2005, 2007). In that framework, a
continuously distributed instrument $Z$ shifts the propensity score $P(Z) =
\mathbb{P}(D = 1 \mid Z)$, and the MTE is identified as $\partial_p E[Y \mid P = p]$.
Our contribution is to show that the overlay design---in which $W_i = D_i Z_i$ and the
treated arm directly reveals $D_i$---provides a natural and operationally distinct
implementation of the local-IV approach: the excluded shifter $S_i$ enters through the
platform's latent propensity $\pi(X_i, S_i)$, and the econometrician observes $P_i$
directly in the arm $Z_i = 1$. Vytlacil (2002) establishes that the latent-index
structure underlying our Assumption~\ref{ass:excluded} is equivalent to the LATE
conditions of Angrist and Imbens (1994), confirming that no additional restrictions
beyond standard IV are invoked.

\textbf{Partial identification.}
Our bounds in Section~\ref{sec:partial} complement the partial identification
literature initiated by Manski (1990). Manski's classic bounds on ATE rely only on
the support of outcomes and can be wide when outcomes are unbounded. Our framework
starts from a stronger baseline---the overlay design identifies $ITT(x)$ and $ATT(x)$,
so the question is only how far $ATE(x)$ can deviate from $ATT(x)$---and derives
sharper bounds by exploiting two economically motivated shape restrictions: monotone
selection on gains (the platform prioritizes high-responders) and nonnegative effects.
The idea of using shape restrictions to tighten partial identification bounds follows
Manski and Pepper (2000).

\textbf{Overlapping and simultaneous platform experiments.}
Tang et al.\ (2010) and Kohavi et al.\ (2020) describe the engineering infrastructure
for overlapping experiments at Google and other large platforms. Their orthogonal
traffic assignment system ensures independent randomization across experiments, which
is the design our paper's Assumption~\ref{ass:overlay} requires for the overlay $Z_i$.
Xiong et al.\ (2020) document the same infrastructure at Tencent. These papers address
the \emph{design} of simultaneous experiments; ours addresses the \emph{identification}
problem that arises when an experimenter cannot enforce orthogonal design and faces an
opaque targeting environment.

\textbf{Combinatorial experiments.}
Ye et al.\ (2025) study the problem of estimating ATEs for all $2^m$ combinations of
$m$ simultaneous A/B tests when only $O(m)$ combinations are observed. Their DeDL
framework uses DML with structured deep neural networks and requires that the treatment
assignment mechanism is known---a natural assumption when the platform controls all
$m$ experiments. Our paper is complementary: we address the case where the assignment
mechanism is unknown, which arises precisely when an external experimenter faces an
opaque platform. The two papers thus address the platform-experiment problem from
opposite sides: Ye et al.\ from inside (platform-controlled, known assignment), ours
from outside (experimenter-controlled, unknown platform propensity).

\textbf{Two-sided platform experiments and interference.}
Johari et al.\ (2022) study bias in A/B tests on two-sided platforms (e.g., ride-sharing,
labor markets) due to market-level interference: treating some users shifts outcomes
for untreated users via market equilibrium effects. Their bias stems from a SUTVA
violation, while ours stems from endogenous selection by an opaque targeting algorithm.
These are distinct but complementary sources of identification failure in platform
settings.

\textbf{Advertising effectiveness.}
Gordon et al.\ (2023) show empirically that non-experimental approaches to advertising
measurement---including DML with observational data---can have errors exceeding 100\%
relative to randomized experiment benchmarks. This motivates the overlay design as
a preferred approach when a clean experiment is infeasible. Our paper provides the
identification theory for the overlay design, explaining precisely what it identifies
(ATT) and under what additional conditions it can recover ATE.

\section{Experimental Setup and Notation}
\label{sec:setup}

Let $i$ index users. The platform makes a latent binary treatment decision
\[
D_i \in \{0,1\},
\]
where $D_i=1$ means that the platform intends to treat unit $i$. The econometrician randomizes a binary overlay
\[
Z_i \in \{0,1\},
\]
where $Z_i=1$ means that the unit is eligible to receive the platform treatment if the platform intends to treat it. Realized treatment receipt is therefore
\[
W_i = D_i Z_i.
\]

The causal object of interest is the effect of the platform treatment $D_i$. Let
\[
Y_i(0), \; Y_i(1)
\]
denote the potential outcomes under no platform treatment and under platform treatment, respectively, and let
\[
\tau_i := Y_i(1)-Y_i(0)
\]
denote the individual treatment effect of the platform treatment. Because the overlay is assumed to act only as a gate, the observed outcome is
\[
Y_i = Y_i(W_i).
\]

The econometrician observes
\[
(Y_i, W_i, Z_i, X_i),
\]
where $X_i$ is a vector of observed covariates. In Section~\ref{sec:pointid}, we additionally introduce an excluded shifter $S_i$.

\begin{assumption}[Randomized Overlay]
\label{ass:overlay}
The overlay assignment $Z_i$ is conditionally independent of the platform's latent intent and potential outcomes:
\[
Z_i \perp (D_i, Y_i(0), Y_i(1)) \mid X_i.
\]
Moreover,
\[
0 < P(Z_i=1 \mid X_i=x) < 1
\]
for all $x$ in the support of $X_i$.
\end{assumption}

\begin{assumption}[One-Sided Receipt]
\label{ass:onesided}
Realized treatment receipt satisfies
\[
W_i = D_i Z_i.
\]
In particular, if $Z_i=0$, then $W_i=0$.
\end{assumption}

Assumption \ref{ass:overlay} states that the econometrician's overlay is randomized within covariate cells. Assumption \ref{ass:onesided} states that the overlay acts only as a gate: a unit receives treatment if and only if both the platform and the econometrician assign treatment.

\section{What a Passive Overlay Identifies}
\label{sec:passive}

For each value $x$ in the support of $X_i$, define the conditional intent-to-treat effect
\[
ITT(x) := E[Y_i \mid Z_i=1, X_i=x] - E[Y_i \mid Z_i=0, X_i=x].
\]
Also define the conditional treatment rate
\[
p(x) := P(W_i=1 \mid Z_i=1, X_i=x).
\]
By Assumption \ref{ass:onesided}, when $Z_i=1$ we have $W_i=D_i$, so
\[
p(x)=P(D_i=1 \mid X_i=x).
\]

The following proposition characterizes the object identified from a passive overlay without any further structure.

\begin{proposition}[Identification of the effect for the platform-treated]
\label{prop:att}
Suppose Assumptions \ref{ass:overlay} and \ref{ass:onesided} hold. Then, for every $x$ such that $p(x)>0$,
\begin{equation}
ITT(x) = E[D_i \tau_i \mid X_i=x],
\label{eq:itt_basic}
\end{equation}
and
\begin{equation}
\frac{ITT(x)}{p(x)} = E[\tau_i \mid D_i=1, X_i=x].
\label{eq:att_basic}
\end{equation}
Hence the passive overlay point-identifies
\[
ATT(x) := E[\tau_i \mid D_i=1, X_i=x].
\]
\end{proposition}

Proposition \ref{prop:att} shows that, in the passive overlay, the econometrician identifies the average causal effect of the platform treatment for the subpopulation that the platform actually selects. In general, this object need not equal the conditional average treatment effect
\[
ATE(x):=E[\tau_i \mid X_i=x].
\]

\section{Point Identification with an Excluded Shifter}
\label{sec:pointid}

This section considers point identification when an excluded shifter is available. Let $S_i$ denote an observed variable that affects the platform's assignment decision but does not directly affect potential outcomes once $X_i$ is held fixed.

\begin{assumption}[Excluded Shifter and Latent Index]
\label{ass:excluded}
There exists a scalar latent variable $U_i$ and a measurable function $\pi(\cdot,\cdot)$ such that:
\begin{enumerate}
    \item The platform assignment rule is
    \[
    D_i = \mathbf{1}\{\pi(X_i,S_i) \ge U_i\}.
    \]
    \item Conditional on $X_i=x$,
    \[
    U_i \sim \mathrm{Unif}[0,1].
    \]
    \item
    \[
    (Y_i(0), Y_i(1), U_i) \perp S_i \mid X_i.
    \]
    \item For each $x$, the support of $\pi(x,S_i)$ conditional on $X_i=x$ contains an interval with nonempty interior.
    \item The overlay remains randomized conditional on $(X_i,S_i)$:
    \[
    Z_i \perp (Y_i(0),Y_i(1),D_i,U_i)\mid (X_i,S_i).
    \]
\end{enumerate}
\end{assumption}

Define the propensity score
\[
P_i := \pi(X_i,S_i),
\]
and define the marginal treatment effect
\[
MTE(x,u) := E[\tau_i \mid X_i=x, U_i=u].
\]

For fixed $(x,p)$, define
\[
\mu_1(x,p) := E[Y_i \mid Z_i=1, X_i=x, P_i=p],
\]
and
\[
\mu_0(x) := E[Y_i \mid Z_i=0, X_i=x].
\]

The next proposition gives the key representation.

\begin{proposition}[Outcome representation]
\label{prop:representation}
Suppose Assumptions \ref{ass:overlay}, \ref{ass:onesided}, and \ref{ass:excluded} hold. Then, for every $(x,p)$ in the support of $(X_i,P_i)$,
\begin{equation}
\mu_1(x,p) = \mu_0(x) + \int_0^p MTE(x,u)\,du.
\label{eq:representation}
\end{equation}
\end{proposition}

\begin{theorem}[Point identification of the MTE and the ATE]
\label{thm:pointid}
Suppose Assumptions \ref{ass:overlay}, \ref{ass:onesided}, and \ref{ass:excluded} hold, and suppose that for each $x$, the function $p \mapsto \mu_1(x,p)$ is differentiable on the interior of the support of $P_i \mid X_i=x$. Then:
\begin{enumerate}
    \item For every $x$ and every $p$ in the interior of the support of $P_i \mid X_i=x$,
    \begin{equation}
    MTE(x,p) = \frac{\partial}{\partial p} \mu_1(x,p).
    \label{eq:mte_derivative}
    \end{equation}
    \item If the support of $P_i \mid X_i=x$ is all of $[0,1]$, then
    \begin{equation}
    ATE(x) := E[\tau_i \mid X_i=x] = \int_0^1 MTE(x,u)\,du.
    \label{eq:ate_conditional}
    \end{equation}
    \item If the support condition in part 2 holds for all $x$, then
    \begin{equation}
    ATE := E[\tau_i] = E\!\left[\int_0^1 MTE(X_i,u)\,du\right].
    \label{eq:ate_population}
    \end{equation}
\end{enumerate}
\end{theorem}

Section~\ref{sec:pointid} shows that when an excluded shifter is available, the overlay environment admits a standard local-IV interpretation. The role of the overlay is operational: because $W_i=D_i$ whenever $Z_i=1$, the econometrician directly observes the platform propensity score in the treated arm.

\section{Partial Identification without an Excluded Shifter}
\label{sec:partial}

We now return to the passive setting in which no excluded shifter is available. In this case the econometrician still observes
\[
p(x)=P(W_i=1 \mid Z_i=1, X_i=x)=P(D_i=1 \mid X_i=x),
\]
but because $p(x)$ is a deterministic function of $X_i$, the variation in treatment propensity cannot be separated from the variation in covariates. A derivative-based local-IV argument is therefore unavailable.

We instead impose a latent-index model and derive sharp bounds.

\begin{assumption}[Passive Latent Index]
\label{ass:passive}
There exists a scalar latent variable $U_i$ such that
\[
D_i = \mathbf{1}\{p(X_i) \ge U_i\},
\]
where, conditional on $X_i=x$,
\[
U_i \sim \mathrm{Unif}[0,1].
\]
\end{assumption}

Under Assumption \ref{ass:passive}, define the marginal treatment effect
\[
MTE(x,u) := E[\tau_i \mid X_i=x, U_i=u].
\]

The following lemma links the observable reduced form to the MTE schedule.

\begin{lemma}[Overlay representation]
\label{lem:overlay}
Suppose Assumptions \ref{ass:overlay}, \ref{ass:onesided}, and \ref{ass:passive} hold. Then, for every $x$,
\begin{equation}
ITT(x) = \int_0^{p(x)} MTE(x,u)\,du.
\label{eq:itt_representation}
\end{equation}
Moreover,
\begin{equation}
ATE(x) := E[\tau_i \mid X_i=x] = \int_0^1 MTE(x,u)\,du.
\label{eq:ate_representation}
\end{equation}
\end{lemma}

To obtain informative bounds, we impose shape restrictions on the MTE schedule.

\begin{assumption}[Monotone Selection on Gains]
\label{ass:monotone}
For every $x$, the function $u \mapsto MTE(x,u)$ is weakly decreasing on $[0,1]$.
\end{assumption}

\begin{assumption}[Nonnegative Effects]
\label{ass:nonnegative}
For every $(x,u)$,
\[
MTE(x,u)\ge 0.
\]
\end{assumption}

Assumption \ref{ass:monotone} says that the platform weakly prioritizes users with larger gains from treatment. This is the econometric translation of the platform's incentive to maximize campaign effectiveness: advertisers and platforms optimize for users with the highest marginal response. Assumption \ref{ass:nonnegative} rules out harmful treatment effects and is used only to obtain the lower bound.

\begin{theorem}[Sharp identification region]
\label{thm:sharp}
Fix $x$ such that $p(x)\in(0,1)$ and $ITT(x)\ge 0$. Suppose Assumptions \ref{ass:overlay}, \ref{ass:onesided}, \ref{ass:passive}, \ref{ass:monotone}, and \ref{ass:nonnegative} hold. Then
\begin{equation}
ITT(x) \le ATE(x) \le \frac{ITT(x)}{p(x)}.
\label{eq:bounds_main}
\end{equation}
Equivalently, since $ATT(x)=ITT(x)/p(x)$,
\begin{equation}
ATE(x) \in \bigl[ITT(x),\, ATT(x)\bigr].
\label{eq:bounds_att}
\end{equation}
The bounds are sharp.
\end{theorem}

\begin{remark}
Under the assumptions of Theorem \ref{thm:sharp}, the platform-treated effect satisfies
\[
ATT(x)=\frac{ITT(x)}{p(x)} \ge ATE(x).
\]
Thus strong selective targeting shows up as a gap between $ATT(x)$ and the lower end of the identified set for $ATE(x)$.
\end{remark}

\subsection{The Selection Ratio}

Theorem~\ref{thm:sharp} suggests a natural summary statistic for the degree of algorithmic
cherry-picking. We define the \emph{Selection Ratio} as:
\begin{equation}
R(x) := \frac{ATT(x)}{ATE(x)}.
\label{eq:selection_ratio}
\end{equation}

The Selection Ratio measures how much the platform's targeted subpopulation outperforms
the population average. If the platform assigns treatments uniformly at random,
$R(x)=1$. If the algorithm perfectly targets the top-$p(x)$ fraction of users by
treatment effect, $R(x)$ can be substantially above 1, indicating strong
cherry-picking.

\begin{corollary}[Bounds on the Selection Ratio]
\label{cor:ratio}
Under the conditions of Theorem~\ref{thm:sharp}, the Selection Ratio is partially identified:
\begin{equation}
R(x) = \frac{ATT(x)}{ATE(x)} \in \left[1,\; \frac{ATT(x)}{ITT(x)}\right] = \left[1,\; \frac{1}{p(x)}\right].
\label{eq:ratio_bounds}
\end{equation}
\end{corollary}

The lower bound $R(x) \ge 1$ follows because $ATE(x) \le ATT(x)$ under monotone
selection. The upper bound $R(x) \le 1/p(x)$ follows because $ATE(x) \ge ITT(x)
= p(x) \cdot ATT(x)$. A lower treatment rate $p(x)$ implies a higher potential
for cherry-picking, as the platform must be highly selective among eligible users.
In practice, advertisers can compare $ATT(x)$ (identifiable from the overlay) with
$ATE(x)$ estimates (from a multi-cell design, if available) to assess how much of
their observed campaign ROI reflects the platform's targeting efficiency rather than
the treatment's intrinsic population value.

\section{Proofs}
\label{sec:proofs}

\subsection*{Proof of Proposition \ref{prop:att}}

When $Z_i=1$, Assumption \ref{ass:onesided} implies $W_i=D_i$, so
\[
Y_i = Y_i(D_i) = Y_i(0) + D_i\bigl(Y_i(1)-Y_i(0)\bigr)
= Y_i(0) + D_i \tau_i.
\]
Therefore, by Assumption \ref{ass:overlay},
\[
E[Y_i \mid Z_i=1, X_i=x]
=
E[Y_i(0) + D_i \tau_i \mid X_i=x]
=
E[Y_i(0)\mid X_i=x] + E[D_i\tau_i \mid X_i=x].
\]
When $Z_i=0$, Assumption \ref{ass:onesided} implies $W_i=0$, so
\[
Y_i = Y_i(0).
\]
Hence
\[
E[Y_i \mid Z_i=0, X_i=x] = E[Y_i(0)\mid X_i=x].
\]
Subtracting yields
\[
ITT(x)=E[D_i\tau_i\mid X_i=x],
\]
which proves \eqref{eq:itt_basic}.

Next, because $D_i$ is binary,
\[
E[D_i\tau_i\mid X_i=x]
=
E[\tau_i \mid D_i=1, X_i=x]P(D_i=1\mid X_i=x).
\]
Since $p(x)=P(D_i=1\mid X_i=x)$, we obtain
\[
E[\tau_i \mid D_i=1, X_i=x]
=
\frac{E[D_i\tau_i\mid X_i=x]}{p(x)}
=
\frac{ITT(x)}{p(x)}.
\]
This proves \eqref{eq:att_basic}. \qed

\subsection*{Proof of Proposition \ref{prop:representation}}

Fix $(x,p)$ in the support of $(X_i,P_i)$. Under Assumptions \ref{ass:onesided} and \ref{ass:excluded}, when $Z_i=1$ and $P_i=p$,
\[
D_i = \mathbf{1}\{p \ge U_i\},
\]
and therefore
\[
Y_i = Y_i(D_i) = Y_i(0) + \mathbf{1}\{U_i\le p\}\tau_i.
\]
Taking conditional expectations,
\[
\mu_1(x,p)
=
E[Y_i(0)\mid Z_i=1, X_i=x, P_i=p]
+
E[\mathbf{1}\{U_i\le p\}\tau_i \mid Z_i=1, X_i=x, P_i=p].
\]
By Assumption \ref{ass:excluded}(5), $Z_i$ is conditionally independent of $(Y_i(0),Y_i(1),U_i)$ given $(X_i,S_i)$. By Assumption \ref{ass:excluded}(3), $(Y_i(0),Y_i(1),U_i)$ is conditionally independent of $S_i$ given $X_i$. Since $P_i=\pi(X_i,S_i)$ is a measurable function of $(X_i,S_i)$, it follows that conditional on $X_i=x$,
\[
(Y_i(0),Y_i(1),U_i) \perp P_i.
\]
Hence
\[
E[Y_i(0)\mid Z_i=1, X_i=x, P_i=p] = E[Y_i(0)\mid X_i=x].
\]
Also,
\[
E[\mathbf{1}\{U_i\le p\}\tau_i \mid Z_i=1, X_i=x, P_i=p]
=
E[\mathbf{1}\{U_i\le p\}\tau_i \mid X_i=x].
\]
Therefore,
\[
\mu_1(x,p)
=
E[Y_i(0)\mid X_i=x]
+
E[\mathbf{1}\{U_i\le p\}\tau_i \mid X_i=x].
\]
Now condition on $U_i$ and use Assumption \ref{ass:excluded}(2):
\[
E[\mathbf{1}\{U_i\le p\}\tau_i \mid X_i=x]
=
\int_0^p E[\tau_i \mid X_i=x, U_i=u]\,du
=
\int_0^p MTE(x,u)\,du.
\]
Finally, when $Z_i=0$, we have $W_i=0$ and hence
\[
\mu_0(x)=E[Y_i\mid Z_i=0,X_i=x]=E[Y_i(0)\mid X_i=x].
\]
Combining the last two displays yields
\[
\mu_1(x,p)=\mu_0(x)+\int_0^p MTE(x,u)\,du.
\]
This proves \eqref{eq:representation}. \qed

\subsection*{Proof of Theorem \ref{thm:pointid}}

By Proposition \ref{prop:representation},
\[
\mu_1(x,p)=\mu_0(x)+\int_0^p MTE(x,u)\,du.
\]
Since $\mu_0(x)$ does not depend on $p$, differentiating both sides with respect to $p$ gives
\[
\frac{\partial}{\partial p}\mu_1(x,p)=MTE(x,p),
\]
which proves \eqref{eq:mte_derivative}.

If the support of $P_i\mid X_i=x$ is all of $[0,1]$, then by Assumption \ref{ass:excluded}(2),
\[
ATE(x)=E[\tau_i\mid X_i=x]
=
\int_0^1 E[\tau_i\mid X_i=x, U_i=u]\,du
=
\int_0^1 MTE(x,u)\,du.
\]
This proves \eqref{eq:ate_conditional}.

Finally,
\[
ATE = E[\tau_i] = E\!\left[E[\tau_i\mid X_i]\right]
= E\!\left[\int_0^1 MTE(X_i,u)\,du\right],
\]
which proves \eqref{eq:ate_population}. \qed

\subsection*{Proof of Lemma \ref{lem:overlay}}

By Proposition \ref{prop:att},
\[
ITT(x)=E[D_i\tau_i\mid X_i=x].
\]
Under Assumption \ref{ass:passive},
\[
D_i=\mathbf{1}\{p(x)\ge U_i\},
\]
so
\[
ITT(x)
=
E[\mathbf{1}\{U_i\le p(x)\}\tau_i \mid X_i=x].
\]
Conditioning on $U_i$ and using $U_i\sim\mathrm{Unif}[0,1]$ conditional on $X_i=x$,
\[
ITT(x)
=
\int_0^{p(x)} E[\tau_i\mid X_i=x,U_i=u]\,du
=
\int_0^{p(x)} MTE(x,u)\,du.
\]
This proves \eqref{eq:itt_representation}.

Similarly,
\[
ATE(x)=E[\tau_i\mid X_i=x]
=
\int_0^1 E[\tau_i\mid X_i=x,U_i=u]\,du
=
\int_0^1 MTE(x,u)\,du,
\]
which proves \eqref{eq:ate_representation}. \qed

\subsection*{Proof of Theorem \ref{thm:sharp}}

By Lemma \ref{lem:overlay},
\[
ITT(x)=\int_0^{p(x)} MTE(x,u)\,du
\quad\text{and}\quad
ATE(x)=\int_0^1 MTE(x,u)\,du.
\]

\paragraph{Upper bound.}
Because $u\mapsto MTE(x,u)$ is weakly decreasing on $[0,1]$, the average value of $MTE(x,u)$ over the treated interval $[0,p(x)]$ is at least as large as the average value over the untreated interval $[p(x),1]$. Formally,
\[
\frac{1}{1-p(x)}\int_{p(x)}^1 MTE(x,u)\,du
\le
\frac{1}{p(x)}\int_0^{p(x)} MTE(x,u)\,du.
\]
Thus
\[
ATU(x):=\frac{1}{1-p(x)}\int_{p(x)}^1 MTE(x,u)\,du
\le
\frac{ITT(x)}{p(x)}=ATT(x).
\]
Using the decomposition
\[
ATE(x)=p(x)ATT(x)+(1-p(x))ATU(x),
\]
we obtain
\[
ATE(x)\le ATT(x)=\frac{ITT(x)}{p(x)}.
\]

\paragraph{Lower bound.}
By Assumption \ref{ass:nonnegative},
\[
\int_{p(x)}^1 MTE(x,u)\,du \ge 0.
\]
Therefore,
\[
ATE(x)
=
\int_0^{p(x)} MTE(x,u)\,du
+
\int_{p(x)}^1 MTE(x,u)\,du
\ge
\int_0^{p(x)} MTE(x,u)\,du
=
ITT(x).
\]

\paragraph{Sharpness.}
Fix $x$, and write
\[
p:=p(x), \qquad I:=ITT(x).
\]
Consider the following two candidate MTE schedules:
\[
m^{L}(u)=
\begin{cases}
I/p, & 0\le u\le p,\\
0, & p<u\le 1,
\end{cases}
\]
and
\[
m^{U}(u)=I/p, \qquad 0\le u\le 1.
\]
Both schedules are weakly decreasing and nonnegative. Moreover,
\[
\int_0^p m^{L}(u)\,du = \int_0^p m^{U}(u)\,du = I.
\]
Hence both schedules generate the same observables $(p(x),ITT(x))$.

Under $m^{L}$,
\[
ATE(x)=\int_0^1 m^{L}(u)\,du = I = ITT(x).
\]
Under $m^{U}$,
\[
ATE(x)=\int_0^1 m^{U}(u)\,du = \frac{I}{p} = ATT(x).
\]
Thus both endpoints are attainable.

Finally, for any $\lambda\in[0,1]$, define
\[
m^{\lambda}(u)=\lambda m^{U}(u)+(1-\lambda)m^{L}(u).
\]
Then $m^{\lambda}(u)$ is again weakly decreasing and nonnegative, satisfies
\[
\int_0^p m^{\lambda}(u)\,du = I,
\]
and yields
\[
\int_0^1 m^{\lambda}(u)\,du
=
\lambda \frac{I}{p} + (1-\lambda)I.
\]
Hence every value in the interval
\[
[ITT(x),ATT(x)]
\]
is attainable under the maintained assumptions. The bounds are therefore sharp. \qed

\subsection*{Proof of Corollary \ref{cor:ratio}}

By Theorem~\ref{thm:sharp}, $ATE(x) \in [ITT(x), ATT(x)]$.

For the lower bound on $R(x)$: $ATE(x) \le ATT(x)$ implies $R(x) = ATT(x)/ATE(x) \ge 1$.

For the upper bound on $R(x)$: $ATE(x) \ge ITT(x) = p(x) \cdot ATT(x)$ implies
$R(x) = ATT(x)/ATE(x) \le ATT(x)/(p(x) \cdot ATT(x)) = 1/p(x)$. \qed

\end{document}
